% INTRODUÇÃO — sem citações, conforme NBR 14724

Os sertões e as caatingas do Ceará e do Rio Grande do Norte abrigam, entre 2004 e 2025, uma das mais extensas expansões de parques eólicos do continente americano. A arquitetura institucional reproduz, sob forma renovável, a racionalidade econômica que orienta a alocação geopolítica de riscos.

Pergunta-se se a expansão eólica no Nordeste brasileiro configura transição energética genuína nos termos da racionalidade ambiental, ou se reproduz, sob discurso verde, lógica análoga à formalizada no início dos anos 1990 sobre territórios de baixa renda.

O objetivo deste artigo consiste em demonstrar que o arranjo fiscal, contratual, territorial, humano e multiespécie do setor eólico cearense e potiguar constitui padrão de captura verde, composto por cinco vetores simultâneos.

Justifica-se a investigação por duas dimensões convergentes. No plano social, comunidades quilombolas, assentadas e pesqueiras têm sido afetadas sem consulta livre, prévia e informada. No plano acadêmico, a literatura jurídica brasileira examina cada vetor de maneira fragmentária, sem amarrar a linearidade institucional que liga a doutrina inicial à expansão eólica nordestina contemporânea.

Metodologicamente, a investigação adota epistemologia dogmática jurídica de tradição analítica, métodos hermenêutico-argumentativo e dedutivo, e técnicas de pesquisa bibliográfica, documental e jurisprudencial. As fontes são primárias e secundárias; terciárias proibidas.

O artigo organiza-se em dois capítulos. O Capítulo 1 constrói a régua conceitual-normativa da transição energética. O Capítulo 2 testa empiricamente essa régua sobre a eólica do Ceará e do Rio Grande do Norte.
